\expandafter\def\csname HistoricalNoteText19\endcsname{Die Berechnung eines Theils dieser Hülfstafel hat Hr. Doctor \textsc{Goldschmidt} ausgeführt.}
\par\clearpage
\SourceLeaf{157}\GaussVTheoryEarlyFolio{153}
\noindent sein. Auch würde man wenig gewinnen, wenn man sämmtliche Glieder, deren Coëfficienten weniger als eine Einheit, oder selbst weniger als 10 Einheiten betragen, unterdrücken wollte, da die Anzahl der übrigen sich doch noch auf 65 belaufen würde. Da nun aber der ganze Werth der Arbeit ungewiss bleiben würde, wenn man sie nicht an einer beträchtlichen Anzahl wirklicher Beobachtungen prüfte, so habe ich die Mühe nicht gescheut, eine Hülfstafel zu berechnen\HistoricalNote{19}, bei deren Gebrauch die Arbeit in hohem Grade abgekürzt und erleichtert, und eben dadurch die Sicherstellung gegen Rechnungsfehler wesentlich befördert wird. Ihre Einrichtung beruhet darauf, dass die Werthe der Componenten in folgende Form gebracht sind
\[
\begin{aligned}
X={}&a^0+a'\cos(\lambda+A')+a''\cos(2\lambda+A'')\\
&+a'''\cos(3\lambda+A''')+a''''\cos(4\lambda+A'''')\\
Y={}&b'\cos(\lambda+B')+b''\cos(2\lambda+B'')\\
&+b'''\cos(3\lambda+B''')+b''''\cos(4\lambda+B'''')\\
Z={}&c^0+c'\cos(\lambda+C')+c''\cos(2\lambda+C'')\\
&+c'''\cos(3\lambda+C''')+c''''\cos(4\lambda+C'''')
\end{aligned}
\]
Die erste Tafel enthält die von $\lambda$ unabhängigen Theile von $X$ und $Z$: in den vier folgenden findet man die Werthe der Hülfswinkel $A'$, $A''$ u. s. w., und der Logarithmen von $a'$, $a''$ u. s. w., alles für die einzelnen Grade der Breite $\varphi=90^0-u$. [Die Tafel ist bei dem vorliegenden Abdruck in einer wegen der Verschiedenheit des Formats etwas abgeänderten Anordnung dem Ende dieser Abhandlung angeschlossen.]

Als Beispiel mag die Rechnung für Göttingen hier Platz finden.

Mit der Breite $+51^0\,32'$ findet man aus den Tafeln
\[
\begin{array}{r|r|r}
a^0=+500{,}8&&c^0=+1465{,}2\\
\log a'=2{,}28980&\log b'=2{,}18900&\log c'=2{,}20204\\
\log a''=1{,}79403&\log b''=2{,}03220&\log c''=2{,}12777\\
\log a'''=1{,}32522&\log b'''=1{,}46845&\log c'''=1{,}43199\\
\log a''''=0{,}59391&\log b''''=0{,}70016&\log c''''=0{,}59091\\
A'=249^0\,30'&B'=358^0\,24'&C'=105^0\,44'\\
A''=311\quad45&B''=64\quad50&C''=165\quad15\\
A'''=234\quad10&B'''=318\quad13&C'''=42\quad22\\
A''''=142\quad26&B''''=232\quad26&C''''=322\quad26
\end{array}
\]
\GaussVTheoryGathering{20}
\par\clearpage
\SourceLeaf{158}\GaussVTheoryEarlyFolio{154}
\noindent und hienach mit der Länge $\lambda=9^0\,56\frac12'$ die Theile von
\[
\begin{array}{r|r|r}
X&Y&Z\\\hline
+500{,}8&&+1465{,}2\\
-35{,}71&+152{,}89&-68{,}99\\
+54{,}76&+9{,}92&-133{,}67\\
-2{,}21&+28{,}77&+8{,}27\\
-3{,}92&+0{,}19&+3{,}90\\\hline
X=+513{,}72&Y=+191{,}77&Z=+1274{,}71
\end{array}
\]
Die weitere Rechnung ergibt dann
\[
\begin{aligned}
\delta&=+20^0\,28'\quad\log\omega=2{,}73907\\
i&=+66\quad43\\
\psi&=1387{,}6\quad\text{oder in der gewöhnlichen Einheit}\\
\psi&=1{,}3876
\end{aligned}
\]

\Article{29}
Die folgende Tafel enthält nun die Vergleichung unsrer Formeln mit den Beobachtungen von 91 Punkten aus allen Theilen der Erde. Da die drei Karten, aus welchen die Data für unsre Rechnung entnommen waren, den Zustand für die neueste Zeit darzustellen bestimmt sind, so wurden auch nur Beobachtungen aus dieser in die Vergleichung aufgenommen, und vorzugsweise von solchen Orten, wo alle drei Elemente des Magnetismus beobachtet sind. Die Forderung einer genauen Gleichzeitigkeit kann jetzt noch nicht gemacht werden, ohne unsern Besitz auf eine äusserst kleine Anzahl herabzusetzen.

Über die hier [am Ende dieses Artikels] zur Vergleichung gebrachten Beobachtungen gebe ich noch folgende Nachweisungen:

Die Intensitätsbestimmungen sind grösstentheils entlehnt aus \textsc{Sabine}'s \textit{Report on the variation of magnetic intensity} (in dem schon oben erwähnten \textit{Seventh Report of the British Association for the advancement of Science}).

Die grosse Anzahl magnetischer Beobachtungen aus dem Russischen Reiche und dem angrenzenden Theile von China verdanken wir

\noindent\hspace*{1em}\textsc{Hansteen} (Poggendorffs Annalen).\par
\noindent\hspace*{1em}\textsc{Erman} (\textit{Reise um die Erde}, und handschriftliche Mittheilungen).\par
\noindent\hspace*{1em}von \textsc{Humboldt} (\textit{Voyage aux regions équinoxiales} T. 13).\par
\par\clearpage
\SourceLeaf{159}\GaussVTheoryEarlyFolio{155}
\noindent\hspace*{1em}\textsc{Fuss} (\textit{Mémoires de l'Académie des Sciences de St. Pétersbourg, Sixième série}).\par
\noindent\hspace*{1em}\textsc{Fedor} (Handschriftlich mitgetheilt durch v. \textsc{Struve}).\par
\noindent\hspace*{1em}\textsc{Reinke} (\textit{Observations météorologiques et magnétiques faites dans l'étendue de l'empire de Russie, rédigées par} A. T. \textsc{Kupffer}, Nr. II).\par
Bei folgenden Örtern wurde das Mittel aus den Bestimmungen mehrerer Beobachter genommen, die zum Theil unter einander grössere Verschiedenheit darbieten. als auf Rechnung der jährlichen Änderungen gesetzt werden kann:

\begin{center}
\begin{minipage}[t]{\linewidth}\centering
(12)\quad\textit{Tobolsk}\par\smallskip
\begin{tabular}{llrr}
Declination.&Hansteen 1828&.\quad.\quad.&$-9^0\,58'$\\
&Erman 1828&.\quad.\quad.\quad.&$-9\quad47$\\
&Fuss 1830&.\quad.\quad.\quad.\quad.&$-11\quad52$\\
&Fedor 1833&.\quad.\quad.\quad.&$-10\quad20$\\
Inclination.&Erman 1828&.\quad.\quad.\quad.\quad.&$71\quad7$\\
&Von Humboldt 1829&.\quad.&$70\quad56$\\
&Fuss 1830&.\quad.\quad.\quad.\quad.\quad.&$71\quad1$\\
&Fedor 1833&.\quad.\quad.\quad.\quad.&$71\quad2$
\end{tabular}
\end{minipage}\par\medskip
\begin{minipage}[t]{\linewidth}\centering
(16)\quad\textit{Catharinenburg}\par\smallskip
\begin{tabular}{llrr}
Declination.&Hansteen 1828&.\quad.\quad.&$-6^0\,27'$\\
&Erman 1828&.\quad.\quad.\quad.&$-7\quad23$\\
&Reinke 1836&.\quad.\quad.\quad.&$-5\quad5$\\
Inclination.&Erman 1828&.\quad.\quad.\quad.\quad.&$69\quad24$\\
&Von Humboldt 1829&.\quad.&$69\quad6$\\
&Fuss 1830&.\quad.\quad.\quad.\quad.&$69\quad19$\\
&Fedor 1832&.\quad.\quad.\quad.\quad.&$69\quad15$
\end{tabular}
\end{minipage}\par\medskip
\begin{minipage}[t]{\linewidth}\centering
(17)\quad\textit{Tomsk}\par\smallskip
\begin{tabular}{llrr}
Declination.&Hansteen 1828&.\quad.\quad.&$-8^0\,32'$\\
&Erman 1829&.\quad.\quad.\quad.&$-8\quad36$\\
Inclination.&Erman 1829&.\quad.\quad.\quad.\quad.&$70\quad59$\\
&Fuss 1830&.\quad.\quad.\quad.\quad.&$70\quad51$
\end{tabular}
\end{minipage}\par\medskip
\begin{minipage}[t]{\linewidth}\centering
(18)\quad\textit{Nishny Nowgorod}\par\smallskip
\begin{tabular}{llrr}
Declination.&Erman 1828&.\quad.\quad.\quad.&$-0^0\,46'$\\
&Fuss 1830&.\quad.\quad.\quad.\quad.&$-0\quad8$
\end{tabular}
\end{minipage}
\end{center}
\GaussVTheoryGathering{20\textsuperscript{*}}
\par\clearpage
\SourceLeaf{160}\GaussVTheoryEarlyFolio{156}
\begin{center}
\begin{minipage}[t]{\linewidth}\centering
(19)\quad\textit{Krasnojarsk}\par\smallskip
\begin{tabular}{llrr}
Declination.&Hansteen 1829&.\quad.\quad.&$-6^0\,43'$\\
&Erman 1829&.\quad.\quad.\quad.&$-6\quad37$\\
&Fedor 1835&.\quad.\quad.\quad.&$-7\quad26$\\
Inclination.&Erman 1829&.\quad.\quad.\quad.\quad.&$70\quad53$\\
&Fedor 1835&.\quad.\quad.\quad.\quad.&$71\quad8$
\end{tabular}
\end{minipage}\par\medskip
\begin{minipage}[t]{\linewidth}\centering
(20)\quad\textit{Kasan}\par\smallskip
\begin{tabular}{llrr}
Inclination.&Erman 1828&.\quad.\quad.\quad.\quad.&$68^0\,21'$\\
&Von Humboldt 1829&.\quad.&$68\quad27$\\
&Fuss 1830&.\quad.\quad.\quad.\quad.&$68\quad26$
\end{tabular}
\end{minipage}\par\medskip
\begin{minipage}[t]{\linewidth}\centering
(21)\quad\textit{Moskwa}\par\smallskip
\begin{tabular}{llrr}
Declination.&Hansteen 1828&.\quad.\quad.&$+3^0\,3'$\\
&Erman 1828&.\quad.\quad.\quad.&$+3\quad1$\\
Inclination.&Erman 1828&.\quad.\quad.\quad.\quad.&$68\quad58$\\
&Von Humboldt 1829&.\quad.&$68\quad57$
\end{tabular}
\end{minipage}\par\medskip
\begin{minipage}[t]{\linewidth}\centering
(30)\quad\textit{Irkuzk}\par\smallskip
\begin{tabular}{llrr}
Declination.&Hansteen 1829&.\quad.\quad.&$-1^0\,37'$\\
&Erman 1829&.\quad.\quad.\quad.&$-1\quad52$\\
&Fuss 1830&.\quad.\quad.\quad.&$-1\quad25$\\
Inclination.&Erman 1829&.\quad.\quad.\quad.\quad.&$68\quad7$\\
&Fuss 1830&.\quad.\quad.\quad.\quad.&$68\quad15$\\
&Fuss 1832&.\quad.\quad.\quad.\quad.&$68\quad20$
\end{tabular}
\end{minipage}\par\medskip
\begin{minipage}[t]{\linewidth}\centering
(36)\quad\textit{Orenburg}\par\smallskip
\begin{tabular}{llrr}
Inclination.&Von Humboldt 1829&.\quad.&$64^0\,41'$\\
&Fedor 1832&.\quad.\quad.\quad.\quad.&$64\quad47$
\end{tabular}
\end{minipage}\par\medskip
\begin{minipage}[t]{\linewidth}\centering
(44)\quad\textit{Troizkosawsk}\par\smallskip
\begin{tabular}{llrr}
Declination.&Hansteen 1829&.\quad.\quad.&$+0^0\,5'$\\
&Erman 1829&.\quad.\quad.&$+0\quad33$\\
&Fuss 1830&.\quad.\quad.\quad.&$-0\quad1$\\
Inclination.&Erman 1829&.\quad.\quad.\quad.\quad.&$66\quad14$\\
&Fuss 1830&.\quad.\quad.\quad.\quad.&$66\quad24$
\end{tabular}
\end{minipage}
\end{center}

\par\clearpage
\SourceLeaf{161}\GaussVTheoryEarlyFolio{157}
Die meisten Bestimmungen in der südlichen Hemisphäre rühren von den Capitaines \textsc{King} und \textsc{Fitz Roy} her, und sind aus einer kleinen Schrift von \textsc{Sabine} (\textit{Magnetic Observations made during the voyages of the ships Adventure and Beagle} 1826—1836) entlehnt.

Die Bestimmungen für die übrigen einzelnen Punkte sind zum Theil auch aus den angeführten Quellen entlehnt; von den andern erwähne ich noch folgende:

(1) Spitzbergen. Beobachter \textsc{Sabine} 1823 (Aus dessen \textit{Account of experiments to determine the figure of the earth}).

(2) Hammerfest. Declination und Inclination im Mittel nach den Bestimmungen von \textsc{Sabine} 1823 (aus angeführtem Werke) und von \textsc{Parry} 1827 (aus dessen \textit{Narrative of an attempt to reach the North Pole}).

(3) Magnetischer Pol, nach \textsc{Ross} 1831 (\textit{Philosophical Transactions} 1834).

(4) Reikiavik nach Beobachtungen von \textsc{Lottin} 1836 (\textit{Voyage en Islande}).

(28) Berlin nach \textsc{Encke} 1836 (\textit{Astronomisches Jahrbuch} 1839).

(38) Göttingen. Die Declination gilt für 1835 Oct. 1 (\textit{Resultate für} 1836 S. 59); die Inclination ist durch Interpolation zwischen von \textsc{Humboldt}s Beobachtung 1826 und \textsc{Forbes} 1837 auf dieselbe Epoche reducirt.

(39) London, nach handschriftlich mitgetheilten Beobachtungen für die Declination von Capitaine \textsc{Ross}; für die Inclination von \textsc{Philips}, \textsc{Fox}, \textsc{Ross}, \textsc{Johnson} und \textsc{Sabine}; die mittlere Epoche für die Declination April 1838, für die Inclination Mai 1838.

(48) Paris für 1835 aus dem \textit{Annuaire} für 1836:

(54) Mailand 1837, von \textsc{Kreil}, nach dessen handschriftlichen Mittheilungen.

(58) Neapel, 1835 nach Beobachtungen von \textsc{Sartorius} und \textsc{Listing}. Die in absolutem Maasse bestimmte Intensität wurde mit dem unten (Art. 31) gegebenen Factor auf die gewöhnliche Einheit reducirt.

(64) Madras 1837 nach \textsc{Taylor}s Beobachtungen, entlehnt aus dem \textit{Journal of the Asiatic Society of Bengal}, Mai 1837.
\par\clearpage
